\begin{theorem}[Theorem 5]
Let $n\ge2$, $p>1$, $0<\tau<n$, let $\mu,\gamma_0,\dots,\gamma_4>0$ with
$\gamma_0\le\gamma_1$, let $d>0$, $M\ge0$, $W\ge0$ and let
$\sigma\colon\R\to\R$, and suppose the Sobolev--Poincar\'e inequality on balls holds in dimension $n$
(\noderef{CitedSobolevPoincareBall}), and assume the cited regularity result
(C1) (\noderef{CitedRegularityEpsPos}), the cited solvability result (C2)
(\noderef{CitedDirichletW1p}) and the cited Dirichlet result (C3) with H\"older
boundary data (\noderef{CitedDirichletHolder}) in the universal forms described
below.  Then there are $N_1,N_2\ge0$, depending only on these
quantities, with the following property.

Let $0\le\varepsilon\le1$, let $\Omega\subseteq\R^n$ be open and let $\Omega'$ be
open, nonempty and bounded and a bounded Lipschitz domain
(\noderef{IsBoundedLipschitzDomain}), with $\overline{\Omega'}\subseteq\Omega$ and
$d\le d(\Omega',\partial\Omega)$ --- ``$\Omega'$ is a bounded Lipschitz domain'' being the hypothesis the paper states at l.~1431, now carried as the predicate \noderef{IsBoundedLipschitzDomain} and no longer only as the three clauses open, nonempty and bounded; those three are kept beside it, being implied by it, so that the proof sites discharging them by name read unchanged ---
and assume the
cited Campanato embedding (C5) on $\Omega'$
(\noderef{CitedCampanatoEmbedding}); assume also the cited H\"older estimates
(C4) (\noderef{CitedHolderOscillationSolution}).  Let $A$ and $B$ satisfy
(Q1) of the thesis--(Q6) of the thesis on $\Omega$ with these parameters and modulus
$\sigma$ (\noderef{EllipticStructureQ12}, \noderef{EllipticStructureQ345},
\noderef{EllipticStructureQ6}).  Let
$u\in W^{1,p}(\Omega)\cap L^\infty(\Omega)$ with weak gradient $g=\nabla u$
(\noderef{MemW1pOn}), $|u|\le M$ a.e.\ and
$\int_\Omega|u|^p\,dx+\int_\Omega|\nabla u|^p\,dx\le W^p$ (that is,
$\|u\|_{W^{1,p}(\Omega)}\le W$), be a weak solution of
(1.1) of the thesis in $\Omega$ (\noderef{IsWeakSolutionOn}).  Then
\begin{itemize}
\item[(i)] $\|\,|\nabla u|\,\|_{L^{p,n-\tau}(\Omega')}\le N_1$, i.e.\
  $|\nabla u|\in L^{p,n-\tau}(\Omega')$ (\noderef{MorreyBoundOn});
\item[(ii)] $\|u\|_{\mathcal L^{\frac n{n-1},\,n+(1-\frac\tau p)(\frac n{n-1})}(\Omega')}\le N_2$
  (\noderef{CampanatoBoundOn});
\item[(iii)] for each $\eta\in[0,1)$, $u$ agrees almost everywhere on $\Omega'$
  with a function in $C^\eta(\overline{\Omega'})$ (\noderef{MemHolderOn}).
\end{itemize}
The norms $N_1,N_2$ are bound after the data, $d(\Omega',\partial\Omega)$, the
function $\sigma$ of (Q6) of the thesis, $\|u\|_{L^\infty(\Omega)}$ and
$\|u\|_{W^{1,p}(\Omega)}$ (represented by $M$ and $W$; no quantity depending on
$|\Omega|$ appears), and before $\varepsilon$,
$\Omega$, $\Omega'$, the operator and $u$: this is the formal content of the
paper's remark on the dependence of these norms.

\medskip
\noindent\textbf{Correction reported.}  Following the paper's Note for the
formalization, $\tau$ ranges over $(0,n)$ and not $[0,n)$: for $\tau=0$ part (i)
would assert $|\nabla u|\in L^{p,n}=L^\infty$, and the proof needs $\tau>0$ to
apply \noderef{CampanatoIterationLemma} with $\beta=n-\tau<\alpha=n$.
\noindent\textbf{On the $W^{1,p}$ datum.}  The hypothesis
$\int_\Omega|u|^p\,dx+\int_\Omega|\nabla u|^p\,dx\le W^p$ is a bound on the
genuine Sobolev norm $\|u\|_{W^{1,p}(\Omega)}$, which is the quantity the
paper's dependency list names.  It must not be replaced by a bound on
$\int_\Omega(1+|\nabla u|^2)^{p/2}dx$: that integrand is pointwise at least $1$,
so such a bound would force $W\ge|\Omega|$, would be unsatisfiable whenever
$|\Omega|=\infty$ --- and the paper states this result for an arbitrary open
$\Omega\subseteq\R^n$, so $\Omega=\R^n$ is an instance --- and would let the
conclusion's norm degrade with $|\Omega|$, which is not on the paper's list.
The weighted integral is still available locally where the proof needs it: since
$(1+t^2)^{p/2}\le2^{p/2}(1+t^p)$ for $t\ge0$,
\[
\int_{B_R}(1+|\nabla u|^2)^{\frac p2}dx
\;\le\;2^{\frac p2}\Bigl(|B_R|+\int_{B_R}|\nabla u|^p\,dx\Bigr)
\;\le\;2^{\frac p2}\bigl(c_nR^{\,n}+W^p\bigr),
\]
in which every term is controlled by $n$, $p$, $R$ and $W$ alone.

The hypothesis \noderef{CitedSobolevPoincareBall} is the cited result (C6) of the
paper's Note for the formalization, carried here as an explicit argument because
the proof uses it --- it enters through the constants $N_1,N_2$ of
\noderef{ComparisonSchemeMorreyCore}, which that node builds with it --- which is what the
Note permits for the results the thesis quotes without proof.  Although it is
bound before the existential for $N_1$ and $N_2$, it introduces no dependence of
$N_1$ and $N_2$ on any quantity outside the list the paper names: (C6) is a statement
about $n$ alone, and $n$ is already bound first, so the hypothesis carries no
numerical content of its own.

\medskip
\noindent\textbf{The three cited results for the frozen operator.}  The cited
regularity result (C1) (\noderef{CitedRegularityEpsPos}), the cited solvability
result (C2) (\noderef{CitedDirichletW1p}) and the cited Dirichlet result (C3)
with H\"older boundary data (\noderef{CitedDirichletHolder}) are assumed here
for \emph{every} field $A'\colon\R^n\to\R^n$ and all parameters
$\varepsilon',\gamma_0',\gamma_1'$ with $\gamma_0'>0$ and $\gamma_1'>0$ --- and,
for (C1), for every open $\Omega_0\subseteq\R^n$ as well --- which are the forms
these three results have in the paper's Note for the formalization, where they
``may be taken as explicit hypotheses \ldots\ where they are used''.

The universal shape is forced rather than convenient.  The operator to which all
three are applied is the frozen operator $h\mapsto A(x_0,u(x_0),h)$, which
varies with the centre $x_0$ and with the value $u(x_0)$; both are chosen inside
the proof, after the existential for $N_1$ and $N_2$, so no single instance of
(C1)--(C3) can be named in the statement.  All three are, in this same
universal form, hypotheses of \noderef{FrozenComparisonDecay}, and the proof
reaches all three through that node: (C1) and (C2) because that node obtains the
decay (4.11) of the thesis from Theorem~1, and (C3) because that node produces the
comparison map.  Like (C6), none of the three introduces a
dependence of $N_1,N_2$ on any quantity outside the list the paper names: each
mentions only $n$ and $p$, both already bound, and each is universally
quantified over its remaining parameters, so none carries numerical content of
its own.  The guards $\gamma_0'>0$, $\gamma_1'>0$ are the thesis' own standing
positivity conditions on the structure constants, and the openness guard on (C1)
is the domain condition \noderef{CitedRegularityEpsPos} is stated for; (C2) and
(C3) need no such guard because \noderef{CitedDirichletW1p} and
\noderef{CitedDirichletHolder} are statements about balls and fix their own
domains.
\textbf{On non-vacuity of the structure hypothesis.}  The statement is
conditioned on \noderef{EllipticStructureQ12}, so it would be vacuously true if
no field satisfied that bundle, which the Note for the formalization forbids at
l.~258.  It does not: \noderef{ModelFieldStructureQ12} exhibits the model field
$A(h)=|h|^{p-2}h$ of (1.3) of the thesis as a member of the bundle at the degenerate
parameter $\varepsilon=0$ on the singular range $1<p<2$, with $\gamma_0=p-1>0$
and $\gamma_1=3-p$.  That is the parameter pair where the question is not
rhetorical: \noderef{EllipticStructureQ12} restricts its differentiability and
its two bounds to $D(\varepsilon)$ precisely because, read on all of $\R^n$, the
conditions as printed in the thesis are satisfied by nothing at all there, and
the cited witness is the machine-checked half of that adjudication.  The citation
is for non-vacuity only; no step of the proof below uses the model field.
\end{theorem}
\begin{proof}
Throughout write
\[
q=\frac n{n-1},\qquad
\theta=n+q\Bigl(1-\frac\tau p\Bigr)
=n+\Bigl(1-\frac\tau p\Bigr)\frac n{n-1},
\qquad
\omega_n=|B_1(0)|\in(0,\infty)
\]
(the Preamble item on the volume of a ball), so that $\theta$ is the Campanato
exponent appearing in (ii) and $q>1$.

The proof follows the thesis' own division of labour at p.~50.  What is
specific to the present setting is the verification that the frozen operator
satisfies (Q1) of the thesis--(Q2) of the thesis at the frozen point, which Step 3 reads
off from (Q1) of the thesis--(Q2) of the thesis and (Q6) of the thesis, and the smallness of
the remainder, which Lemma~5.1 supplies in Step 5.  Everything else is carried
by three nodes that serve this theorem and Theorem~3 alike:
\noderef{InnerDomainRepresentative} produces the inner domain and the
continuous representative (Step 2), \noderef{FrozenComparisonDecay} produces
the comparison map of (4.1) of the thesis and the decay (4.11) of the thesis of its
gradient (Step 4), and \noderef{ComparisonSchemeMorreyCore} carries the whole
passage from (4.11) of the thesis and $(\star)$ to the three conclusions --- the
Morrey decay (4.14) of the thesis, the Morrey bound (i), the Campanato bound (ii)
and the H\"older representative (iii) --- in Step 6.

\emph{Step 1: the constants.}  They are all chosen now, before
$\varepsilon,\Omega,\Omega',A,B,u,g$, which is what the quantifier order of the
statement demands.  Note first that $0<\gamma_1$, since $0<\gamma_0\le\gamma_1$
are hypotheses of the theorem; this is used wherever a cited result asks for a
positive second structure constant.

\begin{enumerate}
\item[(1a)] Put $r=\gamma_1/\gamma_0>0$ and apply
  \noderef{FrozenComparisonDecay} with the parameters $n,p$ and this $r$: its
  hypotheses $2\le n$, $1<p$, $0<r$ hold, its hypothesis
  \noderef{CitedSobolevPoincareBall} is the present theorem's, and its three
  universal hypotheses --- the cited solvability result (C2)
  (\noderef{CitedDirichletW1p}), the cited regularity result (C1)
  (\noderef{CitedRegularityEpsPos}) and the cited Dirichlet result (C3) with
  H\"older boundary data (\noderef{CitedDirichletHolder}) --- are the present
  theorem's (C2), (C1) and (C3) clauses verbatim, in the same universal forms:
  (C2) and (C3) for every field $A'$ and all $\varepsilon',\gamma_0',\gamma_1'$
  with $\gamma_0'>0$ and $\gamma_1'>0$, and (C1) for all of those \emph{and} for
  every open set $\Omega_0$, that clause alone carrying a set in its binder list
  because (C1) is a statement about weak solutions on an open set.  Let $c_0>0$ be the constant
  it produces; $c_0$ depends only on $n,p$ and $\gamma_1/\gamma_0$.
\item[(1b)] Apply \noderef{ComparisonGradientEstimateGeneral} (Lemma~5.1) with
  the parameters $n,p,\mu,\gamma_0,\gamma_1,\gamma_2,\gamma_3,\gamma_4$.  That
  node binds eight numeric guards, plus two cited results, before the two
  clauses of its conclusion, and all ten are hypotheses of the present theorem:
  the guards $2\le n$, $1<p$, $0<\mu$, $0<\gamma_0$, $\gamma_0\le\gamma_1$,
  $0<\gamma_2$, $0<\gamma_3$, $0<\gamma_4$; the Sobolev--Poincar\'e inequality
  on balls in dimension $n$ (\noderef{CitedSobolevPoincareBall}); and the cited
  Dirichlet result (C3) with H\"older boundary data
  (\noderef{CitedDirichletHolder}) for every field $A'$ and all
  $\varepsilon',\gamma_0',\gamma_1'$ with $\gamma_0'>0$ and $\gamma_1'>0$, which
  is the present theorem's (C3) clause verbatim.  Of the two clauses of its
  conclusion take clause (A), which is Lemma 5.1 itself; the present theorem
  assumes no rate for the modulus of (Q6) of the thesis, so clause (B) is not
  invoked.  Read clause (A) at $M$, at $d/2$ in the role of its $d$, and at the
  present $\sigma$, whose guards $0\le M$ and $0<d/2$ hold.
  Let $\sigma_2$ and $R_1>0$ be the function and the radius it produces, so
  $\sigma_2$ is nonnegative, bounded, continuous and increasing with
  $\sigma_2(0)=0$, and both depend only on
  $n,p,\mu,\gamma_0,\gamma_1,\gamma_2,\gamma_3,\gamma_4,M,d$ and $\sigma$.
\item[(1c)] Put
  \[
  R_\ast=\min\Bigl\{R_1,\ \frac d{16}\Bigr\}>0
  \]
  and apply \noderef{ComparisonSchemeMorreyCore} with the parameters $n,p,\tau$,
  with $c_0$ of (1a), with $W$ and $M$, with $d/2$ in the role of its $d$, with
  $R_\ast$ in the role of its $R_1$, and with $\sigma_2$ of (1b): its hypotheses
  $2\le n$, $1<p$, $0<\tau<n$, $0<c_0$, $0\le W$, $0\le M$, $0<d/2$,
  $0<R_\ast$, the four properties of $\sigma_2$ (nonnegativity, continuity,
  monotonicity and $\sigma_2(0)=0$) and the Sobolev--Poincar\'e inequality
  (\noderef{CitedSobolevPoincareBall}) hold by (1a), (1b) and the hypotheses of
  the present theorem.  Let $N_1,N_2\ge0$ be the two constants it produces.
  They depend only on $n,p,\tau,c_0,W,M,d,R_\ast$ and $\sigma_2$, hence --- by
  (1a) and (1b) --- only on
  $n,p,\tau,\mu,\gamma_0,\gamma_1,\gamma_2,\gamma_3,\gamma_4,d,M,W$ and
  $\sigma$, which is the list the statement names, and in particular not on
  $\varepsilon,\Omega,\Omega'$, the data of the problem or $u$ and $g$.
\end{enumerate}

Now fix $\varepsilon\in[0,1]$, the open sets $\Omega,\Omega'$, the fields $A$ and
$B$, the function $u$ and the field $g$ satisfying all the hypotheses.

\emph{Step 2: the inner domain and the continuous representative.}  Apply the
cited H\"older estimates (C4)(b)
(\noderef{CitedHolderOscillationSolution}) with the parameters
$\mu,\gamma_0,\gamma_1,\gamma_2,\gamma_3,\gamma_4,M$: its hypotheses
$0<\gamma_0$, $0<\gamma_1$, $0<\gamma_2$, $0<\gamma_3$, $0<\gamma_4$, $0<\mu$
and $0\le M$ are hypotheses of the present theorem, the second structure
constant being positive by Step 1.  Let $\alpha\in(0,1]$ and $d_0>0$ be the
constants it produces; they depend only on $n,p$, the data and $M$, and in
particular not on $\varepsilon$, which that clause binds after them.  Read its
conclusion at the present $\varepsilon$, whose guards $0\le\varepsilon$ and
$\varepsilon\le1$ are hypotheses of the present theorem.  It applies to the
present data: $\Omega$ is open; $A$ satisfies
(Q1) of the thesis--(Q2) of the thesis on $\Omega$ with the parameters
$p,\varepsilon,\gamma_0,\gamma_1$ (\noderef{EllipticStructureQ12}) and the pair
$A,B$ satisfies (Q3) of the thesis--(Q5) of the thesis with the parameters
$p,\varepsilon,\mu,\gamma_2,\gamma_3,\gamma_4$
(\noderef{EllipticStructureQ345}), both hypotheses of the present theorem; $u$ is
a weak solution of (1.1) of the thesis in $\Omega$ with weak gradient $g$
(\noderef{IsWeakSolutionOn}); and $|u|\le M$ almost everywhere on $\Omega$.  So
for every $y$ and $R_3$ with $\overline B_{R_3}(y)\subseteq\Omega$ and every
$0<\rho\le R_3$,
\[
\operatorname*{osc}\bigl(u,B_\rho(y)\bigr)\le4^\alpha
\max\bigl\{\operatorname*{osc}\bigl(u,B_{R_3}(y)\bigr),\ d_0R_3\bigr\}
\Bigl(\frac\rho{R_3}\Bigr)^\alpha
\tag{1}
\]
(\noderef{EssOscOn}); this is (5.1) of the thesis.

Apply \noderef{InnerDomainRepresentative} with the parameters
$n,p,\alpha,d_0,M,d$ --- its hypotheses $1\le n$, $1\le p$ (because $p>1$),
$0<\alpha$, $\alpha\le1$, $0<d_0$, $0\le M$ and $0<d$ hold --- and with the
data $\Omega,\Omega',u,g$: the set $\Omega$ is open, $\Omega'$ is nonempty and
bounded, $d\le|x-y|$ for every $x\in\Omega'$ and every $y\notin\Omega$,
$u\in W^{1,p}(\Omega)$ with weak gradient $g$ (\noderef{MemW1pOn}),
$|u|\le M$ almost everywhere on $\Omega$, and the oscillation decay is $(1)$.
Let $\Omega_1$ and $U$ be the open set and the function it produces, so that
\begin{itemize}
\item[(2a)] $\Omega_1$ is open and bounded,
  $\overline{\Omega'}\subseteq\Omega_1$ and
  $\overline{\Omega_1}\subseteq\Omega$; in particular
  $\Omega_1\subseteq\Omega$;
\item[(2b)] $\frac d2\le|x-y|$ for every $x\in\Omega'$ and every
  $y\notin\Omega_1$, and $\frac d2\le|x-y|$ for every $x\in\Omega_1$ and every
  $y\notin\Omega$;
\item[(2c)] $\overline B_{d/2}(x)\subseteq\Omega$ for every $x\in\Omega_1$;
\item[(2d)] $U=u$ almost everywhere on $\Omega$, and $U\in W^{1,p}(\Omega)$
  with weak gradient $g$ (\noderef{MemW1pOn});
\item[(2e)] $U$ is continuous on $\Omega_1$, and
  $U\in C^\alpha\bigl(\overline B_{d/16}(x)\bigr)$ (\noderef{MemHolderOn}) for
  every $x\in\Omega_1$.
\end{itemize}

Two consequences of $(2d)$.  First, $|U|\le M$ almost everywhere on $\Omega$,
because $U=u$ almost everywhere there.  Second, $U$ is a weak solution of
(1.1) of the thesis in $\Omega$ with weak gradient $g$ (\noderef{IsWeakSolutionOn}):
apply \noderef{AEModificationInvariance} with the parameters $n,p$ --- its
hypotheses $1\le n$, $1\le p$ hold --- on the open set $\Omega$, with $u$ in the
role of its $f$, $U$ in the role of its $f'$ and the present $g$, the
almost-everywhere equality being that of $(2d)$; its clause (4), taken at the
present fields $A$ and $B$, transports the weak-solution property of $u$ to $U$.

\emph{Step 3: the frozen operator.}  Fix $x_1\in\Omega_1$ and
$0<R\le R_\ast$; note $R\le R_\ast\le d/16$, so
$\overline B_R(x_1)\subseteq\overline B_{d/2}(x_1)\subseteq\Omega$ by
$(2c)$.  Put
\[
z_1=U(x_1),\qquad A^0(h)=A(x_1,z_1,h),\qquad
\widetilde A=A^0-A^0(0),
\]
so that $\widetilde A(0)=0$ by construction.  Apply
\noderef{FrozenSolutionStructureQ12} with the present
$n,p,\varepsilon,\gamma_0,\gamma_1$, the open set $\Omega$, the field $A$, the
modulus $\sigma$, the point $x_1$ and the value $z_1$.  Its hypotheses are
$2\le n$, $1<p$, $0\le\varepsilon$, $0<\gamma_0$ and $0<\gamma_1$ --- the last
from $0<\gamma_0\le\gamma_1$, all fixed in Step 1 --- the openness of $\Omega$,
the membership $x_1\in\Omega$, which holds because $x_1\in\Omega_1$ and
$\overline\Omega_1\subseteq\Omega$ by $(2a)$, and the two bundles
(Q1) of the thesis--(Q2) of the thesis (\noderef{EllipticStructureQ12}) and
(Q6) of the thesis (\noderef{EllipticStructureQ6}), both hypotheses of the present
theorem.  It gives \emph{both}
\[
\mathrm{EllipticStructureQ12}\;n\;p\;\varepsilon\;\gamma_0\;\gamma_1\;
\mathrm{univ}\;\bigl((x,z,h)\mapsto A^0(h)\bigr)
\quad\text{and}\quad
\mathrm{EllipticStructureQ12}\;n\;p\;\varepsilon\;\gamma_0\;\gamma_1\;
\mathrm{univ}\;\bigl((x,z,h)\mapsto\widetilde A(h)\bigr),
\]
and it is the \emph{normalised} one that Step 4 consumes.

\emph{Why this is a citation and no longer an inline derivation.}  An earlier
version of this proof discharged all six clauses of
\noderef{EllipticStructureQ12} by hand here, on the ground that the twin
Theorem~Theorem~3 of the thesis's dedicated node for the frozen operator is stated for the
frozen \emph{gradient} $\nabla_hf^0$ of a variational integrand and so has
nothing to consume on the solution side, which carries the field $A$ directly.
That remains true of the twin's node and is now beside the point: the
equation-side analogue exists, takes exactly the field $A$ together with the two
bundles named above, and performs the normalisation internally.  The one step
that cannot be taken naively is inside it --- (Q1) of the thesis and (Q2) of the thesis
hold only for \emph{almost every} $x$, whereas $x_1$ is a single point and hence
null, so they say nothing at the centre until they are upgraded to hold at every
point of $\Omega$, and that upgrade is what consumes (Q6) of the thesis (pm-0034).
The passage from $A^0$ to $\widetilde A$ is the paper's own normalisation at
l.~542, ``since $\operatorname{div}A(0)=0$ we may assume without loss of
generality that $A(0)=0$'', and the two Carath\'eodory clauses ride along under
it unchanged, subtracting a constant vector altering neither measurability in
$x$ nor continuity in $(z,h)$.

\emph{Step 4: the comparison map and the decay (4.11) of the thesis.}  Apply the
second part of \noderef{FrozenComparisonDecay} --- the same instance as in
(1a), so the constant it produces is the $c_0$ fixed there --- with the data
$\varepsilon$, $\alpha$, $\gamma_0$, $\gamma_1$, the field $A^0$, the centre
$x_1$, the radius $R$, the function $U$ and the field $g$.  Its hypotheses
hold:
\begin{itemize}
\item $0\le\varepsilon$ and $\varepsilon\le1$; $0<\alpha$ and $\alpha\le1$
  (Step 2); $0<\gamma_0$ and $0<\gamma_1$ (Step 1); $\gamma_1\le r\gamma_0$,
  with equality by the definition of $r$ in (1a); and $0<R$;
\item $\widetilde A=A^0-A^0(0)$ satisfies (Q1) of the thesis--(Q2) of the thesis on
  $\R^n$ with the parameters $p,\varepsilon,\gamma_0,\gamma_1$
  (\noderef{EllipticStructureQ12}), which is Step 3;
\item $U\in W^{1,p}(B_R(x_1))$ with weak gradient $g$ (\noderef{MemW1pOn}), by
  $(2d)$ and clause (1) of \noderef{WeakGradientRestrict} applied with $q=p$ to
  the open subset $B_R(x_1)\subseteq\Omega$: its outer hypotheses are $1\le n$
  (from $2\le n$), $1\le p$ (from $1<p$), the openness of $\Omega$ and of
  $B_R(x_1)$, and the inclusion $B_R(x_1)\subseteq\Omega$, which holds because
  $B_R(x_1)\subseteq\overline B_R(x_1)\subseteq\Omega$ by Step 3, and its
  clause-(1) hypothesis is the $U\in W^{1,p}(\Omega)$ of $(2d)$;
\item $U\in C^\alpha\bigl(\overline B_R(x_1)\bigr)$ (\noderef{MemHolderOn}):
  by $(2e)$ at the centre $x_1\in\Omega_1$ we have
  $U\in C^\alpha(\overline B_{d/16}(x_1))$, and
  $\overline B_R(x_1)\subseteq\overline B_{d/16}(x_1)$ because
  $R\le R_\ast\le d/16$; a function of class $C^\alpha$ on a set is of class
  $C^\alpha$ on every subset of it, the continuity restricting and the same
  nonnegative constant serving.
\end{itemize}
It produces functions $v,V$ and a vector field $g_v$ such that the bundle
$\mathrm{FrozenComparisonOn}(n,p,x_1,R,A^0,U,U,g,v,V,g_v)$ holds
(\noderef{FrozenComparisonOn}) --- so in particular $v$ is a weak solution of
$\operatorname{div}A^0(\nabla v)=0$ in $B_R(x_1)$
(\noderef{IsWeakSolutionOn}), $v\in W^{1,p}(B_R(x_1))$ with weak gradient $g_v$
(\noderef{MemW1pOn}), and $v-U\in W^{1,p}_0(B_R(x_1))$ with weak gradient
$g_v-g$ (\noderef{MemW1pZeroOn}) --- and such that
\[
\int_{B_\rho(x_1)}\bigl(\varepsilon+|g_v|^2\bigr)^{\frac p2}dx
\le c_0\Bigl(\frac\rho R\Bigr)^{n}
\int_{B_R(x_1)}\bigl(\varepsilon+|g_v|^2\bigr)^{\frac p2}dx
\qquad(0<\rho\le R),
\tag{4.11}
\]
which is the first of the two estimates the comparison scheme asks for.

\emph{Step 5: the smallness $(\star)$.}  Apply
\noderef{ComparisonGradientEstimateGeneral} (Lemma~5.1) with the parameters of
(1b) --- so clause (A) of its conclusion, with $d/2$ in the role of its
$d$, and with the $\sigma_2$ and $R_1$ it produced --- and with the data $\varepsilon$, $\Omega$,
$\Omega_1$ in the role of its $\Omega'$, the fields $A$ and $B$,
$U$ in the role of its $u$ \emph{and} of its $U$, the field $g$, the centre
$x_1$, the radius $R$, and the $v,V,g_v$ of Step 4.  Its hypotheses hold:
$0\le\varepsilon\le1$; $\Omega$ is open; $\Omega_1$ is open and
bounded with $\overline{\Omega_1}\subseteq\Omega$ and, by $(2b)$,
$d/2\le|x-y|$ for $x\in\Omega_1$, $y\notin\Omega$; the three
structure bundles (Q1) of the thesis--(Q2) of the thesis
(\noderef{EllipticStructureQ12}), (Q3) of the thesis--(Q5) of the thesis
(\noderef{EllipticStructureQ345}) and (Q6) of the thesis
(\noderef{EllipticStructureQ6}) hold on $\Omega$ with the stated parameters;
the cited H\"older estimates (C4) are assumed
(\noderef{CitedHolderOscillationSolution}); $U\in W^{1,p}(\Omega)$ with weak
gradient $g$, $|U|\le M$ almost everywhere on $\Omega$, and $U$ is a
weak solution of (1.1) of the thesis in $\Omega$ with weak gradient $g$, all three by
Step 2; $x_1\in\Omega_1$; $0<R\le R_\ast\le R_1$;
$\overline B_R(x_1)\subseteq\Omega$ (Step 3); and the frozen comparison bundle
of Step 4, whose frozen operator is $h\mapsto A(x_1,U(x_1),h)=A^0(h)$, as that
node requires.  Its conclusion is
\[
\int_{B_R(x_1)}|g-g_v|^p\,dx
\le\sigma_2(R)\int_{B_R(x_1)}(1+|g|^2)^{\frac p2}dx .
\tag{$\star$}
\]

\emph{Step 6: the comparison scheme, and the three conclusions.}  Apply
\noderef{ComparisonSchemeMorreyCore} with the parameters of (1c) --- so with
$n,p,\tau$, with $c_0$ of (1a), with $W$ and $M$, with $d/2$ in the role of its
$d$, with $R_\ast$ in the role of its $R_1$ and with $\sigma_2$ of (1b) --- and
with the data $\varepsilon$, $\Omega$, $\Omega_1$ in the role of
its $\Omega''$, $\Omega'$, the function $u$ and the field $g$.  Its hypotheses
hold:
\begin{itemize}
\item $0\le\varepsilon\le1$; $\Omega$, $\Omega_1$ and $\Omega'$ are
  open, $\Omega'$ is nonempty by hypothesis and $\Omega_1$ is bounded, and
  $\overline{\Omega'}\subseteq\Omega_1$ and
  $\overline{\Omega_1}\subseteq\Omega$, all by $(2a)$;
\item the two separations, both with the constant $d/2$, which are $(2b)$;
\item $\Omega'$ a bounded Lipschitz domain
  (\noderef{IsBoundedLipschitzDomain}), a hypothesis of the present theorem,
  which that node now carries because the (C5) instance it assumes requires it;
\item the cited Campanato embedding (C5) on $\Omega'$ at the exponent
  $n/(n-1)$ (\noderef{CitedCampanatoEmbedding}) and the Sobolev--Poincar\'e
  inequality on balls (\noderef{CitedSobolevPoincareBall}) are hypotheses of the
  present theorem;
\item $u\in W^{1,p}(\Omega)$ with weak gradient $g$ (\noderef{MemW1pOn}),
  $|u|\le M$ almost everywhere on $\Omega$, and
  $\int_\Omega|u|^p\,dx+\int_\Omega|\nabla u|^p\,dx\le W^p$: these are
  hypotheses of the present theorem;
\item the comparison bundle, which is where Steps 2--5 are consumed: the
  centre $x_1\in\Omega_1$ and the radius $0<R\le R_\ast$ were arbitrary in
  Steps 3--5, so for every such pair the field $g_v$ produced in
  Step 4 satisfies (4.11) of the thesis by Step 4 and $(\star)$ by Step 5 --- and
  $\overline B_R(x_1)\subseteq\Omega$ holds automatically for such a pair, by
  $(2c)$ and $R_\ast\le d/16$, so the node's guard
  $\overline B_R(x_1)\subseteq\Omega$ is met.  Note that $R_\ast\le R_1$,
  which is what Step 5 requires, and $R_\ast\le d/16$, which is what Steps 3
  and 4 require.
\end{itemize}
Its three conclusions, with the $N_1$ and $N_2$ of (1c), are verbatim the three
assertions of the present theorem: the Morrey bound
$\|\,|\nabla u|\,\|_{L^{p,n-\tau}(\Omega')}\le N_1$
(\noderef{MorreyBoundOn}), which is (i); the Campanato bound
$\|u\|_{\mathcal L^{q,\theta}(\Omega')}\le N_2$ with
$\theta=n+(1-\frac\tau p)\frac n{n-1}$ (\noderef{CampanatoBoundOn}), which is
(ii); and, for each $\eta\in[0,1)$, a function $w$ equal to $u$ almost
everywhere on $\Omega'$ and of class $C^\eta(\overline{\Omega'})$
(\noderef{MemHolderOn}), which is (iii).  This completes the proof.
\end{proof}
